
\chapter{Construcción de la gramática producto de una gramática de concatenación de rango y un lenguaje regular finito}
\label{cap:interseccion}

La literatura provee construcciones de la intersección con un lenguaje
regular para las gramáticas libres del contexto y las de
concatenación de rango simple~\cite{barhillel1961, bertsch2001rcg}. En las gramáticas
de concatenación de rango, una misma variable puede aparecer más de una
vez en el lado derecho de una cláusula. Esta no linealidad introduce una
dificultad que las construcciones anteriores no contemplan. Este capítulo aporta una
construcción que, a partir de una gramática de concatenación de rango
$G$ y un autómata finito acíclico $A$, produce una nueva gramática de
concatenación de rango $G'$, la \textbf{gramática
producto}\footnote{El término se usa por analogía con el autómata
producto, la construcción estándar para la intersección de dos
lenguajes regulares~\cite{hopcroft2006introduction}.},
cuyo lenguaje es exactamente $L(G) \cap L(A)$.

La construcción se apoya en una noción de rango adaptada al autómata,
que se introduce en la sección siguiente.

\section{El rango sobre un autómata}

En una RCG el reconocimiento se hace por rangos; en un autómata, por
estados. Para reconocer una cadena con ambos hay que coordinarlos: de
cada rango se necesita saber en qué estado queda el autómata al empezar
y al terminar de leer su subcadena. Pero el rango $(i, j)$ es un par de
posiciones en una cadena fija, ajeno a los estados del autómata.

Hace falta, entonces, una noción de rango ligada a los estados del
autómata. Para ello, cada subcadena se caracteriza por el par de estados
entre los que transita el autómata al leerla. Para fijar ese par hace
falta saber en qué estado termina el autómata al comenzar en un estado y
leer una cadena, lo que da lugar a la siguiente definición.

\begin{definicion}
Sea $A = (Q, \Sigma, \delta, q_0, F)$ un autómata finito determinista. La
función de transición $\delta$ se extiende a cadenas mediante la
\textbf{función de transición extendida}
$\hat\delta : Q \times \Sigma^* \to Q$, que se define por
\[
  \hat\delta(q, \varepsilon) = q \quad\text{y}\quad
  \hat\delta(q, wa) = \delta(\hat\delta(q, w), a)
\]
para $w \in \Sigma^*$ y $a \in \Sigma$~\cite{hopcroft2006introduction}.
Así, $\hat\delta(q, u)$ es el estado en el que el autómata termina al
comenzar en el estado $q$ y leer la cadena $u$.
\end{definicion}

En los ejemplos de este capítulo se retoma la modelación del
capítulo~\ref{cap:modelacion-sat} de la fórmula
$x_1 \wedge (x_2 \vee x_1)$, cuya secuencia de ocurrencias es
$x_1\,x_2\,x_1$. Su autómata booleano $A$ se muestra en la
figura~\ref{fig:automata-booleano}. Es finito y acíclico: reconoce las
cadenas de longitud tres que hacen verdadera la fórmula cuando cada
ocurrencia se trata por separado. Su estado inicial es $q_0$, su único
estado de aceptación es $V$, y su lenguaje es $L(A) = \{101, 110, 111\}$.

Por ejemplo, al leer $101$ sobre $A$ se extiende $\hat\delta$ símbolo a
símbolo,
\[
  \hat\delta(q_0, 1) = q_1, \quad \hat\delta(q_0, 10) = q_2 \quad\text{y}\quad
  \hat\delta(q_0, 101) = V,
\]
un estado de aceptación; para la cadena $100$, en cambio, se obtiene
$\hat\delta(q_0, 100) = F$. La figura~\ref{fig:delta-ext-ej} muestra
estos cuatro caminos en $A$.

\begin{figure}[H]
\centering
\setlength{\tabcolsep}{10pt}
\begin{tabular}{cc}
\begin{tikzpicture}[->, >=stealth, shorten >=1pt, auto, semithick, scale=0.46, transform shape]
  \node[state, initial, fill=blue!30] (q0) at (0,0) {$q_0$};
  \node[state, fill=red!30] (q1) at (2.6,1.3) {$q_1$};
  \node[state] (q4) at (2.6,-1.7) {$q_4$};
  \node[state] (q3) at (5.2,2.3) {$q_3$};
  \node[state] (q2) at (5.2,0.3) {$q_2$};
  \node[state] (q5) at (5.2,-1.7) {$q_5$};
  \node[state, accepting] (V) at (7.8,1.3) {$V$};
  \node[state] (F) at (7.8,-1.7) {$F$};
  \path
    (q0) edge [orange, very thick] node {$1$} (q1)
    (q0) edge node[swap] {$0$} (q4)
    (q1) edge node {$1$} (q3)
    (q1) edge node {$0$} (q2)
    (q3) edge node {$0,1$} (V)
    (q2) edge node {$1$} (V)
    (q2) edge node[swap] {$0$} (F)
    (q4) edge node {$0,1$} (q5)
    (q5) edge node {$0,1$} (F);
\end{tikzpicture} &
\begin{tikzpicture}[->, >=stealth, shorten >=1pt, auto, semithick, scale=0.46, transform shape]
  \node[state, initial, fill=blue!30] (q0) at (0,0) {$q_0$};
  \node[state] (q1) at (2.6,1.3) {$q_1$};
  \node[state] (q4) at (2.6,-1.7) {$q_4$};
  \node[state] (q3) at (5.2,2.3) {$q_3$};
  \node[state, fill=red!30] (q2) at (5.2,0.3) {$q_2$};
  \node[state] (q5) at (5.2,-1.7) {$q_5$};
  \node[state, accepting] (V) at (7.8,1.3) {$V$};
  \node[state] (F) at (7.8,-1.7) {$F$};
  \path
    (q0) edge [orange, very thick] node {$1$} (q1)
    (q0) edge node[swap] {$0$} (q4)
    (q1) edge node {$1$} (q3)
    (q1) edge [orange, very thick] node {$0$} (q2)
    (q3) edge node {$0,1$} (V)
    (q2) edge node {$1$} (V)
    (q2) edge node[swap] {$0$} (F)
    (q4) edge node {$0,1$} (q5)
    (q5) edge node {$0,1$} (F);
\end{tikzpicture} \\[0.4em]
{\small (a) $\hat\delta(q_0, 1) = q_1$} &
{\small (b) $\hat\delta(q_0, 10) = q_2$} \\[1.4em]
\begin{tikzpicture}[->, >=stealth, shorten >=1pt, auto, semithick, scale=0.46, transform shape]
  \node[state, initial, fill=blue!30] (q0) at (0,0) {$q_0$};
  \node[state] (q1) at (2.6,1.3) {$q_1$};
  \node[state] (q4) at (2.6,-1.7) {$q_4$};
  \node[state] (q3) at (5.2,2.3) {$q_3$};
  \node[state] (q2) at (5.2,0.3) {$q_2$};
  \node[state] (q5) at (5.2,-1.7) {$q_5$};
  \node[state, accepting, fill=red!30] (V) at (7.8,1.3) {$V$};
  \node[state] (F) at (7.8,-1.7) {$F$};
  \path
    (q0) edge [orange, very thick] node {$1$} (q1)
    (q0) edge node[swap] {$0$} (q4)
    (q1) edge node {$1$} (q3)
    (q1) edge [orange, very thick] node {$0$} (q2)
    (q3) edge node {$0,1$} (V)
    (q2) edge [orange, very thick] node {$1$} (V)
    (q2) edge node[swap] {$0$} (F)
    (q4) edge node {$0,1$} (q5)
    (q5) edge node {$0,1$} (F);
\end{tikzpicture} &
\begin{tikzpicture}[->, >=stealth, shorten >=1pt, auto, semithick, scale=0.46, transform shape]
  \node[state, initial, fill=blue!30] (q0) at (0,0) {$q_0$};
  \node[state] (q1) at (2.6,1.3) {$q_1$};
  \node[state] (q4) at (2.6,-1.7) {$q_4$};
  \node[state] (q3) at (5.2,2.3) {$q_3$};
  \node[state] (q2) at (5.2,0.3) {$q_2$};
  \node[state] (q5) at (5.2,-1.7) {$q_5$};
  \node[state, accepting] (V) at (7.8,1.3) {$V$};
  \node[state, fill=red!30] (F) at (7.8,-1.7) {$F$};
  \path
    (q0) edge [orange, very thick] node {$1$} (q1)
    (q0) edge node[swap] {$0$} (q4)
    (q1) edge node {$1$} (q3)
    (q1) edge [orange, very thick] node {$0$} (q2)
    (q3) edge node {$0,1$} (V)
    (q2) edge node {$1$} (V)
    (q2) edge [orange, very thick] node[swap] {$0$} (F)
    (q4) edge node {$0,1$} (q5)
    (q5) edge node {$0,1$} (F);
\end{tikzpicture} \\[0.4em]
{\small (c) $\hat\delta(q_0, 101) = V$} &
{\small (d) $\hat\delta(q_0, 100) = F$}
\end{tabular}
\caption[Caminos de los ejemplos de $\hat\delta$]{Caminos en $A$ de los cuatro ejemplos de $\hat\delta$. En cada
uno, el estado de partida $q_0$ va en azul y el de llegada en rojo;
las transiciones resaltadas forman el camino, a lo largo del cual se lee
la cadena correspondiente.}
\label{fig:delta-ext-ej}
\end{figure}

La función $\hat\delta$ determina, para cada cadena y cada estado de
partida, el estado en que el autómata termina, lo que permite definir el
rango sobre el autómata como un par de estados y precisar cuándo una
cadena lo realiza.

\begin{definicion}
Un \textbf{rango sobre $A$} es un par de estados $(q, q') \in Q \times Q$.
Una cadena $u$ \textbf{realiza} el rango $(q, q')$ si $\hat\delta(q, u) = q'$.
\end{definicion}

Un rango sobre $A$ describe a la vez todas las cadenas que lo realizan.
Cuando a una variable se le asocia el rango $(q, q')$, la gramática debe
derivar la subcadena correspondiente, que a la vez debe realizar ese
rango.

En el ejemplo anterior, al leer $101$ se recorre en $A$ el camino
$q_0\,1\,q_1\,0\,q_2\,1\,V$\footnote{Recuérdese que en la notación de
camino del capítulo~\ref{cap:vacio-interseccion-rcg-regular} estados y
símbolos se alternan.}. Así, el primer símbolo realiza
$(q_0, q_1)$; el segundo, $(q_1, q_2)$; el tercero, $(q_2, V)$; la
subcadena $10$ realiza $(q_0, q_2)$; y la cadena entera $101$ realiza
$(q_0, V)$, que parte del estado inicial y termina en el de aceptación.

En la sección siguiente, las nociones planteadas para las cadenas se
generalizan al caso de los argumentos.

\section{El rango de un argumento sobre un autómata}\label{sec:rango-argumento}

La sección anterior estableció cuándo una cadena realiza un rango sobre
$A$. Un argumento de una cláusula es una secuencia de terminales y
variables, y las variables no son símbolos que el autómata pueda leer.
Hace falta entonces decir qué rango sobre $A$ le corresponde a un
argumento de esa forma.

\begin{definicion}
Sea $\alpha = \alpha_1 \cdots \alpha_n$ un argumento, donde cada
$\alpha_i \in T \cup V$. Una \textbf{asignación de rangos sobre $A$} para
$\alpha$ es una función $\rho$ que asigna a cada $\alpha_i$ un rango
sobre $A$, $\rho(\alpha_i) = (q_i, q'_i)$, con dos condiciones:
\begin{enumerate}
  \item $q'_i = q_{i+1}$ para $1 \le i < n$;
  \item si $\alpha_i$ es un terminal $a$, entonces $\delta(q_i, a) = q'_i$.
\end{enumerate}
Cuando existe una asignación así, el argumento $\alpha$ realiza
el rango sobre $A$ $(q_1, q'_n)$.
\end{definicion}

Por ejemplo, considérese el argumento $W_0\,1\,W_1\,1\,W_2$ del lado
izquierdo de la cláusula
$Eq_{x_1}(W_0\,1\,W_1\,1\,W_2) \to L_0(W_0)\,L_1(W_1)\,L_0(W_2)$ de
$G_{\mathrm{cons}}$, que reconoce las cadenas de longitud tres con $1$ en
las dos posiciones de $x_1$. Sobre la cadena consistente $101$, con
$W_0 = \varepsilon$, $W_1 = 0$ y $W_2 = \varepsilon$, la asignación de
rangos sobre $A$ es la siguiente, con $W_0$ y $W_2$ vacíos en los
extremos:

\begin{center}
\begin{tikzpicture}[->, >=stealth, shorten >=1pt, auto, node distance=2.2cm, semithick]
  \node[state] (s0) {$q_0$};
  \node[state] (s1) [right=of s0] {$q_1$};
  \node[state] (s2) [right=of s1] {$q_2$};
  \node[state] (s3) [right=of s2] {$V$};
  \path (s0) edge              node {$1$} (s1);
  \path (s1) edge [dashed]     node {$W_1$} (s2);
  \path (s2) edge              node {$1$} (s3);
\end{tikzpicture}
\end{center}

Esta asignación cumple las dos condiciones de la definición de
asignación de rangos sobre $A$. La primera, que el final de cada rango
sea el inicio del siguiente, se ve al recorrer los rangos: $W_0$ realiza
$(q_0, q_0)$; el primer $1$, $(q_0, q_1)$; $W_1$, $(q_1, q_2)$; el
segundo $1$, $(q_2, V)$; y $W_2$, $(V, V)$. La segunda, que cada terminal
$a$ con rango $(q_i, q'_i)$ cumpla $\delta(q_i, a) = q'_i$, la satisfacen
los dos terminales $1$: $\delta(q_0, 1) = q_1$ y $\delta(q_2, 1) = V$. El
argumento realiza entonces $(q_0, V)$, que parte de $q_0$ y termina en el
estado de aceptación $V$.

Con la cláusula del símbolo $0$, $Eq_{x_1}(W_0\,0\,W_1\,0\,W_2)$, se da
el caso contrario. Al leer la cadena consistente $000$ se recorre el
camino $q_0\,0\,q_4\,0\,q_5\,0\,F$, de modo que el argumento realiza
$(q_0, F)$, que no termina en el estado de aceptación: $000$ es
consistente, pero hace falsa la fórmula.

La sección siguiente presenta la construcción de la gramática producto.

\section{Construcción de la gramática producto}

En esta sección se desarrolla la construcción de la gramática
producto. Primero se presenta el algoritmo. A continuación se impone
una restricción sobre la gramática de entrada. Después se establece
una condición para que las ocurrencias de una misma variable reciban
el mismo rango. Luego se da el pseudocódigo. Por último se demuestra
su correctitud y se analizan su tamaño y la complejidad de construirla.

\subsection{Algoritmo}\label{subsec:algoritmo}

El objetivo es producir una gramática $G' = (N', T, V, P', S')$ que
reconozca las cadenas que $G$ reconoce y $A$ acepta. Una cadena $w$
pertenece a $L(A)$ si el
rango sobre $A$ que $w$ realiza parte del estado inicial y termina en
un estado final, así que $G'$ no solo debe reconocer las cadenas de
$L(G)$, sino reflejar también qué rango sobre $A$ realiza cada
argumento de cada predicado durante la derivación. Para reflejarlo,
los predicados de $G'$ se diferencian según el rango sobre $A$ que
realiza cada uno de sus argumentos: cada predicado de $G$ da lugar a
un predicado de $G'$ por cada vector de rangos sobre $A$.

\begin{definicion}
Sean $B$ un predicado de $G$ de aridad $k$ y
$\vec q = ((q_1, q'_1), \ldots, (q_k, q'_k))$ un vector de $k$ rangos
sobre $A$, uno por cada argumento de $B$. El par $(B, \vec q)$ se
llama \textbf{predicado indexado} y se denota $B[\vec q]$.
\end{definicion}

Por ejemplo, el autómata $A$ tiene ocho estados, así que cada predicado
de $G_{\mathrm{cons}}$ da lugar a hasta sesenta y cuatro predicados
indexados, uno por cada par de estados. En el reconocimiento de $101$
aparecen $S[(q_0, V)]$,
$Eq_{x_1}[(q_0, V)]$ y $Eq_{x_2}[(q_0, V)]$, junto con predicados
indexados de $L_0$ y $L_1$.

El predicado inicial de $G'$ es un predicado $S'$ de aridad 1. El
conjunto de predicados $N'$ consta de los predicados indexados y del
predicado inicial $S'$.

Las cláusulas de $G'$ se generan a partir de las cláusulas de $G$ y de
las asignaciones de rangos sobre $A$ a sus ocurrencias de símbolos. Para
cada cláusula
\[
  B_0(\alpha_{0,1}, \ldots, \alpha_{0,k_0}) \to
  B_1(\alpha_{1,1}, \ldots, \alpha_{1,k_1}) \cdots
  B_m(\alpha_{m,1}, \ldots, \alpha_{m,k_m})
\]
de $G$ y cada asignación que cumpla las condiciones de la
sección~\ref{sec:rango-argumento}, indexando cada predicado $B_j$ por el
vector $\vec q_j$ de los rangos de sus argumentos se obtiene una
\textbf{cláusula indexada}\footnote{Una cláusula indexada es una
cláusula de la gramática producto $G'$, formada por predicados
indexados.}:
\[
  B_0[\vec q_0](\alpha_{0,1}, \ldots, \alpha_{0,k_0}) \to
  B_1[\vec q_1](\alpha_{1,1}, \ldots, \alpha_{1,k_1}) \cdots
  B_m[\vec q_m](\alpha_{m,1}, \ldots, \alpha_{m,k_m}).
\]

Por ejemplo, en la cláusula $L_1(1X) \to L_0(X)$ de $G_{\mathrm{cons}}$,
con el terminal $1$ en el rango $(q_0, q_1)$, pues $\delta(q_0, 1) = q_1$,
y la variable $X$ en $(q_1, q_1)$, el argumento $1X$ realiza $(q_0, q_1)$. La
cláusula indexada que se obtiene es
\[
  L_1[(q_0, q_1)](1X) \to L_0[(q_1, q_1)](X).
\]
La construcción genera una cláusula indexada por cada asignación de
rangos que cumpla las condiciones, de modo que cada cláusula de
$G_{\mathrm{cons}}$ da lugar a muchas.

Cuando el lado derecho de una cláusula de $G$ es vacío, solo se indexa el lado izquierdo: la cláusula $B(\alpha_1, \ldots, \alpha_k) \to \varepsilon$
produce las cláusulas indexadas $B[\vec q](\alpha_1, \ldots, \alpha_k) \to
\varepsilon$, una por cada asignación de rangos a sus argumentos. El
argumento $\varepsilon$ realiza solo el rango $(q, q)$ para cada estado
$q$, así que la cláusula $L_0(\varepsilon) \to \varepsilon$ produce ocho
cláusulas en $G'$, una por cada estado del autómata:
\begin{align*}
  L_0[(q_0, q_0)](\varepsilon) &\to \varepsilon, &
  L_0[(q_1, q_1)](\varepsilon) &\to \varepsilon, \\
  L_0[(q_2, q_2)](\varepsilon) &\to \varepsilon, &
  L_0[(V, V)](\varepsilon) &\to \varepsilon,
\end{align*}
y de igual forma para los estados $q_3$, $q_4$, $q_5$ y $F$.

El conjunto de cláusulas $P'$ consta de las cláusulas anteriores y de
las del predicado inicial $S'$, que dependen de la condición de
aceptación de $A$: una cadena se acepta cuando el rango que realiza
parte del estado inicial $q_0$ y termina en un estado final. Por eso la
construcción solo usa los predicados indexados $S[(q_0, q_F)]$ con
$q_F \in F$, y, para cada uno, $G'$ tiene una cláusula
\[
  S'(X) \to S[(q_0, q_F)](X),
\]
donde $S$ es el predicado inicial de $G$. Una cadena $w$ pertenece a
$L(G')$ si $S'(w)$ se deriva a la cadena vacía.

En el ejemplo de $G_{\mathrm{cons}}$ sobre $A$, $F = \{V\}$\footnote{Aquí
$F$ es el conjunto de estados de aceptación del autómata, no el estado
$F$ del autómata booleano.}, y el rango
$(q_0, V)$ parte del estado inicial y termina en el estado de
aceptación. El predicado inicial $S'$ tiene la cláusula:
\[
  S'(X) \to S[(q_0, V)](X).
\]

La construcción conserva todas las cláusulas que produce, tanto las
que provienen de las cláusulas de $G$ como las del predicado inicial
$S'$. Las que participan en el reconocimiento de una cadena se
determinan durante la derivación, no durante la construcción.

La construcción produce los predicados $N'$, las cláusulas $P'$ y el
predicado inicial $S'$, y mantiene los terminales $T$ y las variables
$V$ de $G$. Los cinco componentes completan la tupla
$G' = (N', T, V, P', S')$, con lo que la construcción de la gramática
producto termina.

Para ilustrar el reconocimiento en la gramática producto $G'$, se
describe la derivación de la cadena
$101 \in L(G_{\mathrm{cons}}) \cap L(A)$, el testigo de que la fórmula
es satisfacible. Cada predicado se indexa por el rango que su argumento
realiza en la lectura de $101$ por $A$, cuyo camino es
$q_0\,1\,q_1\,0\,q_2\,1\,V$:
\begin{align*}
  S'(101) &\to S[(q_0, V)](101), \\
  S[(q_0, V)](101) &\to Eq_{x_1}[(q_0, V)](101)\; Eq_{x_2}[(q_0, V)](101), \\
  Eq_{x_1}[(q_0, V)](101) &\to L_0[(q_0, q_0)](\varepsilon)\; L_1[(q_1, q_2)](0)\; L_0[(V, V)](\varepsilon), \\
  Eq_{x_2}[(q_0, V)](101) &\to L_1[(q_0, q_1)](1)\; L_1[(q_2, V)](1), \\
  L_1[(q_1, q_2)](0) &\to L_0[(q_2, q_2)](\varepsilon), \\
  L_1[(q_0, q_1)](1) &\to L_0[(q_1, q_1)](\varepsilon), \\
  L_1[(q_2, V)](1) &\to L_0[(V, V)](\varepsilon),
\end{align*}
y cada predicado indexado $L_0[(q, q)](\varepsilon)$ deriva en la cadena
vacía por la cláusula terminal. La figura~\ref{fig:derivacion-101}
muestra el árbol de derivación, con cada predicado indexado por el rango
que realiza su argumento. El primer paso aplica la cláusula del
predicado inicial $S'$; el segundo, la cláusula no lineal de $S$, que
asocia el mismo rango $(q_0, V)$ a las dos ocurrencias de $X$. Como la
derivación de $S'(101)$ termina en la cadena vacía, $G'$ reconoce $101$.

La derivación de $101$ usa solo unos pocos de los predicados y cláusulas
de $G'$. La construcción genera muchos más que no aparecen en el
reconocimiento de ninguna cadena de $L(G_{\mathrm{cons}}) \cap L(A)$:
los predicados indexados cuyo rango no parte de $q_0$ ni termina en $V$,
como $Eq_{x_1}[(q_0, F)]$, quedan sin uso, junto con las cláusulas que
los tienen en el lado izquierdo.

\begin{figure}[H]
\centering
\begin{forest}
for tree={font=\footnotesize, parent anchor=south, child anchor=north,
          l sep=7mm, s sep=3mm, inner sep=2pt}
[{$S'(101)$}
  [{$S[(q_0, V)](101)$}
    [{$Eq_{x_1}[(q_0, V)](101)$}
      [{$L_0[(q_0, q_0)](\varepsilon)$} [{$\varepsilon$}]]
      [{$L_1[(q_1, q_2)](0)$} [{$L_0[(q_2, q_2)](\varepsilon)$} [{$\varepsilon$}]]]
      [{$L_0[(V, V)](\varepsilon)$} [{$\varepsilon$}]]
    ]
    [{$Eq_{x_2}[(q_0, V)](101)$}
      [{$L_1[(q_0, q_1)](1)$} [{$L_0[(q_1, q_1)](\varepsilon)$} [{$\varepsilon$}]]]
      [{$L_1[(q_2, V)](1)$} [{$L_0[(V, V)](\varepsilon)$} [{$\varepsilon$}]]]
    ]
  ]
]
\end{forest}
\caption[Árbol de derivación de $101$ en $G'$]{Árbol de derivación de $101$ en la gramática producto $G'$.
Cada predicado va indexado por el rango que realiza su argumento; las
hojas $\varepsilon$ corresponden a las cláusulas terminales
$L_0[(q, q)](\varepsilon) \to \varepsilon$.}
\label{fig:derivacion-101}
\end{figure}

Por conveniencia, la construcción anterior se restringe a una clase
de gramáticas de entrada que evita casos aparte. La sección siguiente
la precisa.

\subsection{Restricción sobre la gramática de entrada}

Cada predicado se indexa por los rangos sobre $A$ de sus argumentos.
El lado derecho de una cláusula fija los rangos de las variables que
contiene. Si una variable apareciera en el lado izquierdo pero no en el
lado derecho, el lado derecho no fijaría su rango, y habría que tratarla como un
caso aparte.

Para evitarlo, las gramáticas de entrada se restringen a las que
cumplen una propiedad sobre sus cláusulas: en cada cláusula, toda
variable del lado izquierdo aparece también en el lado derecho.

\begin{definicion}
Una gramática de concatenación de rango es \textbf{no borrante}
si en cada cláusula toda variable del lado izquierdo aparece también en el
lado derecho.
\end{definicion}

Se pide además que los argumentos del lado derecho sean variables, de modo
que cada variable del lado derecho sea un argumento por sí sola.

\begin{definicion}
Una gramática de concatenación de rango es \textbf{no combinatoria} si
en cada cláusula los argumentos del lado derecho son variables.
\end{definicion}

Por ejemplo, las cláusulas de $G_{\mathrm{cons}}$ cumplen ambas: en
$Eq_{x_1}(W_0\,1\,W_1\,1\,W_2) \to L_0(W_0)\,L_1(W_1)\,L_0(W_2)$ las
variables $W_0, W_1, W_2$ del lado izquierdo aparecen en el lado derecho,
cada una como un argumento por sí sola, y $L_0(\varepsilon) \to
\varepsilon$ no tiene variables. Una cláusula como $A(XY) \to B(X)$ no
cumple la condición de no borrante, pues $Y$ figura en el lado izquierdo
y falta en el lado derecho, y $A(X) \to B(X1)$ no cumple la de no
combinatoria, pues el argumento del lado derecho $X1$ no es una variable.

Exigir ambas propiedades no descarta ningún lenguaje: toda gramática de
concatenación de rango admite una equivalente que las
cumple~\cite{boullier1998proposal, kallmeyer2010parsing}.

Además de la restricción sobre la gramática de entrada, hace falta una
condición sobre los rangos que se asignan, que se detalla a
continuación.

\subsection{Variables repetidas en una cláusula}

En las gramáticas de concatenación de rango, una misma variable puede
aparecer varias veces en una cláusula. En la construcción de la
sección~\ref{subsec:algoritmo}, cada argumento del lado derecho se indexa por
separado,
sin que se fije el mismo rango para las ocurrencias de una misma
variable, que representan una sola cadena.
Sin una condición que lo imponga, se les pueden asignar rangos
distintos, y entonces $G'$ reconocería cadenas que no están en la
intersección.

Por eso es necesario añadir la siguiente condición: en cada cláusula,
a todas las ocurrencias de una variable se les asigna el mismo rango.
Más formalmente, si una variable $X$ aparece en dos argumentos de la
cláusula con asignaciones de rangos $\rho$ y $\rho'$, entonces
$\rho(X) = \rho'(X)$.

Por ejemplo, la cláusula no lineal de $G_{\mathrm{cons}}$,
\[
  S(X) \to Eq_{x_1}(X)\, Eq_{x_2}(X),
\]
tiene la variable $X$ en tres argumentos: el del lado izquierdo y los
dos del lado derecho. Sobre $A$, la cadena $000$ es consistente, pero
hace falsa la fórmula: realiza $(q_0, F)$, que no termina en un estado
de aceptación.

Sin esta condición existe la cláusula
\[
  S[(q_0, V)](X) \to Eq_{x_1}[(q_0, F)](X)\, Eq_{x_2}[(q_0, F)](X),
\]
con el predicado del lado izquierdo indexado en $(q_0, V)$ y los del
lado derecho en $(q_0, F)$. Como $000$ realiza $(q_0, F)$ y es
consistente, los dos predicados $Eq$ del lado derecho se derivan en la
cadena vacía; con la cláusula $S'(X) \to S[(q_0, V)](X)$, la gramática
$G'$ reconocería $000$, que está fuera de la intersección.

Con la condición, las tres ocurrencias de $X$ llevan el mismo rango, así
que la única cláusula cuyo lado izquierdo se indexa en $(q_0, V)$ es
\[
  S[(q_0, V)](X) \to Eq_{x_1}[(q_0, V)](X)\, Eq_{x_2}[(q_0, V)](X).
\]
Como $000$ realiza $(q_0, F)$ y no $(q_0, V)$, el predicado
$Eq_{x_1}[(q_0, V)](000)$ no se deriva en la cadena vacía, así que $000$
queda fuera del lenguaje de $G'$, como debería ser.

La repetición de variables es justo lo que separa a las RCG de las sRCG,
y la condición no hace falta para estas últimas: en una sRCG
cada variable aparece una sola vez en el lado derecho, así que no hay varias
ocurrencias que coordinar. Imponerla es lo que permite tratar RCG y no
solo sRCG.

Ya se presentaron la idea del algoritmo, la restricción sobre la
gramática de entrada y la condición sobre las ocurrencias de una misma
variable; con ellas se fija la construcción. A continuación se describe
el pseudocódigo.

\subsection{Pseudocódigo}

El procedimiento $\proc{Product-Grammar}$ recibe una gramática de
concatenación de rango no combinatoria y
no borrante $G = (N, T, V, P, S)$ y un autómata finito
determinista acíclico $A = (Q, \Sigma, \delta, q_0, F)$; devuelve
$G' = (N', T, V, P', S')$ con $L(G') = L(G) \cap L(A)$.

\begin{codebox}
\Procname{$\proc{Product-Grammar}(G, A)$}
\li $P' = \emptyset$
\li \For each clause $c = B_0(\alpha_{0,1}, \ldots, \alpha_{0,k_0}) \to {}$
\zi \qquad $B_1(\alpha_{1,1}, \ldots, \alpha_{1,k_1}) \cdots B_m(\alpha_{m,1}, \ldots, \alpha_{m,k_m})$ in $P$
\li \Do
        let $\sigma_1, \sigma_2, \ldots$ be the symbol occurrences of $c$ and its empty arguments
\li     $\ell = $ number of these
\li     \For each $\mu = ((q_1, q'_1), \ldots, (q_\ell, q'_\ell)) \in (Q \times Q)^\ell$
\li     \Do
            \If $\proc{Consistent}(c, \mu)$
\li         \Then
                \For $j = 0$ \To $m$
\li             \Do
                    $\vec{q}_j = (\proc{Range}(\alpha_{j,1}, \mu), \ldots, \proc{Range}(\alpha_{j,k_j}, \mu))$
                \End
\li             add $B_0[\vec{q}_0](\alpha_{0,1}, \ldots, \alpha_{0,k_0}) \to {}$
\zi \qquad\qquad $B_1[\vec{q}_1](\alpha_{1,1}, \ldots, \alpha_{1,k_1}) \cdots B_m[\vec{q}_m](\alpha_{m,1}, \ldots, \alpha_{m,k_m})$ to $P'$
            \End
        \End
    \End
\li \For each $q \in F$
\li \Do
        add $S'(X) \to S[(q_0, q)](X)$ to $P'$
    \End
\li $N' = \{\, S' \,\} \cup \{\, B' : B' \text{ occurs in } P' \,\}$
\li \Return $(N', T, V, P', S')$
\end{codebox}

En la línea~2 se recorren las cláusulas de $G$. Para cada una, en las
líneas~3 y~4 se cuentan los símbolos que aparecen en ella, con
repeticiones, y sus argumentos vacíos. En la línea~5 se recorren todas
las posibles asignaciones de rangos. De la línea~6 a la~9, con $\proc{Consistent}$ se conservan
solo las asignaciones en las que cada rango termina donde empieza el
siguiente, cada terminal $a$ con rango $(q, q')$ cumple
$\delta(q, a) = q'$, y todas las ocurrencias de una variable reciben el
mismo rango; por cada una se indexa cada predicado por los rangos de sus
argumentos con $\proc{Range}$, y se agrega la cláusula a $P'$. En las líneas~10 y~11 se
agregan las cláusulas del predicado inicial $S'$, una por cada estado
final. En la línea~12 se agregan a $N'$ el
predicado $S'$ y los predicados que aparecen en $P'$.

En $\proc{Consistent}$ se comprueban las condiciones sobre $\mu$:

\begin{codebox}
\Procname{$\proc{Consistent}(c, \mu)$}
\li \For each non-empty argument $\alpha = \alpha_1 \cdots \alpha_n$ of $c$, with $\mu(\alpha_i) = (q_i, q'_i)$
\li \Do
        \For $i = 1$ \To $n$
\li     \Do
            \If $i < n$ \kw{and} $q'_i \neq q_{i+1}$
\li         \Then \Return \const{false}
            \End
\li         \If $\alpha_i$ is a terminal $a$ \kw{and} $\delta(q_i, a) \neq q'_i$
\li         \Then \Return \const{false}
            \End
        \End
    \End
\li \For each empty argument $\alpha$ of $c$, with $\mu(\alpha) = (q, q')$
\li \Do
        \If $q \neq q'$
\li     \Then \Return \const{false}
        \End
    \End
\li \For each variable $X$ of $c$
\li \Do
        \If two occurrences of $X$ have different ranges under $\mu$
\li     \Then \Return \const{false}
        \End
    \End
\li \Return \const{true}
\end{codebox}

En las líneas~1--9 de $\proc{Consistent}$ se comprueban las condiciones
sobre los rangos de los argumentos; en las líneas~10--12, que las
ocurrencias de una variable reciban el mismo rango.

$\proc{Range}(\alpha, \mu)$ es el rango que el argumento $\alpha$
realiza bajo $\mu$:

\begin{codebox}
\Procname{$\proc{Range}(\alpha, \mu)$}
\li \If $\alpha = \varepsilon$
\li \Then \Return $\mu(\alpha)$
    \End
\li let $\alpha = \alpha_1 \cdots \alpha_n$, with $\mu(\alpha_i) = (q_i, q'_i)$
\li \Return $(q_1, q'_n)$
\end{codebox}

Para $\alpha = \alpha_1 \cdots \alpha_n$ es $(q_1, q'_n)$, el inicio del
primer símbolo y el final del último; para $\alpha = \varepsilon$ es el
rango $(q, q)$ que $\mu$ le asigna, el rango que $\alpha$ realiza según
la sección~\ref{sec:rango-argumento}.

En la próxima subsección se demuestra que $G'$ reconoce
exactamente $L(G) \cap L(A)$.

\subsection{Correctitud}

En toda la sección $G$ es no combinatoria y no borrante. La igualdad
$L(G') = L(G) \cap L(A)$ se prueba por las dos inclusiones.
La clave es una propiedad de los índices: una cláusula indexada solo
reconoce cadenas que realizan los rangos con que se indexan sus
argumentos. De ahí, $A$ acepta toda cadena que $G'$ reconoce; y, al
borrar los índices, también la reconoce $G$. En el otro sentido, se
parte de una derivación de $G$ y se indexa cada predicado con los rangos
sobre $A$ que sus argumentos realizan al leer la cadena; el resultado es
una derivación de $G'$.

Sea $B_0[\vec q_0](u_1, \ldots, u_{k_0})$ el predicado indexado $B_0[\vec q_0]$
aplicado a las cadenas $u_1, \ldots, u_{k_0}$, una por cada argumento, con
$\vec q_0 = ((q_1, q'_1), \ldots, (q_{k_0}, q'_{k_0}))$.

\begin{lema}\label{lema:proyeccion}
Si $B_0[\vec q_0](u_1, \ldots, u_{k_0})$ se deriva a $\varepsilon$ en $G'$,
entonces $B_0(u_1, \ldots, u_{k_0})$ se deriva a $\varepsilon$ en $G$.
\end{lema}

\begin{proof}
Las cláusulas que intervienen en una derivación de
$B_0[\vec q_0](u_1, \ldots, u_{k_0})$ son cláusulas indexadas, cada una
proveniente de una cláusula de $P$ a la que solo se le agregan los
índices de los predicados; el predicado $S'$ no aparece, pues no figura
en el lado derecho de ninguna cláusula. Al borrar los índices, cada cláusula
indexada se reduce a su cláusula de $P$, y se obtiene una derivación de
$B_0(u_1, \ldots, u_{k_0})$ en $G$.
\end{proof}

La hipótesis de que $G$ es no combinatoria es clave en el segundo lema.

\begin{lema}\label{lema:realizacion}
Si $B_0[\vec q_0](u_1, \ldots, u_{k_0})$ se deriva a $\varepsilon$ en $G'$,
entonces cada $u_l$ realiza el rango $(q_l, q'_l)$.
\end{lema}

\begin{proof}
La prueba es por inducción en la cantidad de pasos de la derivación de
$B_0[\vec q_0](u_1, \ldots, u_{k_0})$ a $\varepsilon$. La hipótesis de inducción
es que el lema vale para toda derivación con menos pasos.

La derivación comienza con una cláusula de $P'$ de lado izquierdo
$B_0[\vec q_0]$, que proviene de una cláusula de $P$,
\[
  c = B_0(\alpha_{0,1}, \ldots, \alpha_{0,k_0}) \to B_1(\alpha_{1,1}, \ldots, \alpha_{1,k_1}) \cdots B_m(\alpha_{m,1}, \ldots, \alpha_{m,k_m}),
\]
y de una asignación $\mu$ de un rango a cada ocurrencia de símbolo de $c$.
Recuérdese que esta asignación cumple tres condiciones. La primera es que
en cada argumento el final del rango de un símbolo sea el inicio del
siguiente. La segunda
es que cada terminal $a$ con rango $(q, q')$ cumpla $\delta(q, a) = q'$.
La tercera es que todas las ocurrencias de una variable reciban el mismo
rango. La cláusula indexada que resulta es
\[
  B_0[\vec q_0](\alpha_{0,1}, \ldots, \alpha_{0,k_0}) \to B_1[\vec q_1](\alpha_{1,1}, \ldots, \alpha_{1,k_1}) \cdots B_m[\vec q_m](\alpha_{m,1}, \ldots, \alpha_{m,k_m}),
\]
con $(q_l, q'_l)$ el rango que el argumento $\alpha_{0,l}$ del lado
izquierdo realiza bajo $\mu$. La cláusula se instancia con una sustitución $\theta$ de las
variables por subcadenas, de modo que $u_l = \theta(\alpha_{0,l})$.

En el caso base, el lado derecho de $c$ es vacío ($m = 0$), de modo que
\[
  c = B_0(\alpha_{0,1}, \ldots, \alpha_{0,k_0}) \to \varepsilon,
\]
y la derivación tiene un solo paso. Como el lado derecho es vacío y $G$
es no borrante, el lado izquierdo no tiene variables, pues toda variable
suya tendría que aparecer en el lado derecho. Así, cada argumento
$\alpha_{0,l}$ es una cadena de terminales, y por la primera y la segunda
condición $u_l = \alpha_{0,l}$ realiza $(q_l, q'_l)$.

En el paso inductivo, el lado derecho no es vacío, y cada predicado
\[
  B_j[\vec q_j](\theta(\alpha_{j,1}), \ldots, \theta(\alpha_{j,k_j}))
\]
se deriva a $\varepsilon$ en menos pasos. Como $G$ es no combinatoria, cada
argumento del lado derecho es una sola variable $X$, con un único rango
$\mu(X)$ en todas sus ocurrencias por la tercera condición, y el rango
que ese argumento realiza es $\mu(X)$. Por la hipótesis de inducción, $\theta(X)$ realiza $\mu(X)$
para toda variable $X$ del lado derecho; como $G$ es no borrante, lo
mismo vale para toda variable de $c$.

Sea entonces $\alpha_{0,l}$ un argumento del lado izquierdo. Por la
segunda condición, cada terminal de $\alpha_{0,l}$ realiza el rango que
$\mu$ le asigna; y para cada variable $X$, $\theta(X)$ realiza $\mu(X)$,
que por la tercera condición es el rango de todas sus ocurrencias. Como
el final de cada rango es el inicio del siguiente, por la primera
condición, la concatenación de estas subcadenas,
$u_l = \theta(\alpha_{0,l})$, realiza la concatenación de sus rangos,
$(q_l, q'_l)$. Si
$\alpha_{0,l} = \varepsilon$, su rango $(q_l, q'_l)$ tiene $q_l = q'_l$,
por ser el de un argumento vacío, y $u_l = \varepsilon$ lo realiza.
\end{proof}

Con los dos lemas se obtiene el teorema.

\begin{teorema}\label{teo:correctitud}
$L(G') = L(G) \cap L(A)$.
\end{teorema}

\begin{proof}
Inclusión $L(G') \subseteq L(G) \cap L(A)$. Sea $w \in L(G')$;
entonces $S'(w)$ se deriva a $\varepsilon$. Las únicas cláusulas de
lado izquierdo $S'$ son $S'(X) \to S[(q_0, q_F)](X)$ con $q_F \in F$, así que el
primer paso es $S'(w) \to S[(q_0, q_F)](w)$ para algún $q_F \in F$, y
$S[(q_0, q_F)](w)$ se deriva a $\varepsilon$. Por el
lema~\ref{lema:proyeccion}, $S(w)$ se deriva a $\varepsilon$ en $G$, esto
es, $w \in L(G)$. Por el lema~\ref{lema:realizacion}, $w$ realiza
$(q_0, q_F)$, es decir, $\hat\delta(q_0, w) = q_F$; como $q_F \in F$,
$w \in L(A)$. Luego $w \in L(G) \cap L(A)$.

Inclusión $L(G) \cap L(A) \subseteq L(G')$. Sea
$w = a_1 \cdots a_n \in L(G) \cap L(A)$. El camino de la lectura de $w$
por $A$ es $p_0\,a_1\,p_1 \cdots a_n\,p_n$, con $p_0 = q_0$ y
$p_i = \delta(p_{i-1}, a_i)$; como $w \in L(A)$, $p_n \in F$. La subcadena
en las posiciones $i$ a $j$ realiza el rango sobre $A$ $(p_i, p_j)$.

Como $w \in L(G)$, hay una derivación de $S(w)$ a $\varepsilon$ en $G$,
en la que cada predicado se aplica a subcadenas de $w$ en posiciones
fijas. Al indexar cada predicado por los rangos sobre $A$ de sus
argumentos se obtiene una derivación en $G'$, porque cada cláusula
indexada pertenece a $P'$. En efecto, sea
\[
  B_0(\alpha_{0,1}, \ldots, \alpha_{0,k_0}) \to B_1(\alpha_{1,1}, \ldots, \alpha_{1,k_1}) \cdots B_m(\alpha_{m,1}, \ldots, \alpha_{m,k_m})
\]
una cláusula instanciada en la
derivación, y $\mu$ la asignación que a cada ocurrencia de símbolo le
asigna el rango sobre $A$ de la posición que ocupa. La asignación $\mu$
cumple las tres condiciones: los símbolos contiguos de un argumento
ocupan posiciones contiguas, así que el final de cada rango es el inicio
del siguiente, y un argumento vacío en la posición $i$ ocupa el rango
$(p_i, p_i)$; para un terminal $a$ en la posición $i{+}1$ se
tiene $\delta(p_i, a) = p_{i+1}$; y las ocurrencias de una
variable ocupan las mismas posiciones, luego reciben el mismo rango. Por tanto, la construcción produce la cláusula
indexada, que pertenece a $P'$.

El lado izquierdo $S(w)$ se indexa por el rango $(p_0, p_n) =
(q_0, p_n)$. Como $p_n \in F$, la cláusula
$S'(X) \to S[(q_0, p_n)](X)$ pertenece a $P'$, así que $S'(w)$ se
deriva a $\varepsilon$ y $w \in L(G')$.
\end{proof}

La gramática producto es correcta, pero su tamaño puede ser exponencial
en la entrada, como se ve a continuación.

\subsection{Tamaño y complejidad}

El tamaño de $G'$ y el costo de construirla dependen del número de
predicados $|N|$ y cláusulas $|P|$ de $G$, de los $|Q|$ estados de $A$,
de la aridad máxima $a$ y del máximo $\ell$ de ocurrencias de símbolo y
argumentos vacíos por cláusula.

Cada predicado de aridad $k$ se indexa con un vector de $k$ rangos sobre
$A$. Un rango es un par de estados, de modo que hay $|Q|^2$ rangos y
$\bigl(|Q|^2\bigr)^k = |Q|^{2k}$ vectores posibles. Como la aridad máxima
es $a$, por cada uno de los $|N|$ predicados hay a lo sumo $|Q|^{2a}$
predicados indexados. Con el predicado inicial,
\[
  |N'| = O(|N|\,|Q|^{2a}).
\]
Para cada cláusula con $\ell_c$ ocurrencias de símbolo se recorren, en la
línea~5 del procedimiento $\proc{Product-Grammar}$, las $|Q|^{2\ell_c}$
asignaciones de un rango a cada una, y por cada una que
cumple las tres condiciones se agrega a lo sumo una cláusula indexada;
las del predicado inicial son una por estado final. Así,
\[
  |P'| = O(|P|\,|Q|^{2\ell} + |F|) = O(|P|\,|Q|^{2\ell}).
\]
Como los terminales son los de $G$, el tamaño de $G'$ es
\[
  |G'| = |N'| + |T| + |P'| = O\bigl(|N|\,|Q|^{2a} + |P|\,|Q|^{2\ell} + |T|\bigr).
\]
El costo de construir $G'$ es el número de asignaciones recorridas por el
$O(\ell_c)$ de verificar y procesar cada una:
\[
  O(|P|\,\ell\,|Q|^{2\ell}).
\]
Los factores $|Q|^{2a}$ y $|Q|^{2\ell}$ son exponenciales en la aridad
máxima $a$ y en $\ell$. Por eso el tamaño de $G'$ es exponencial cuando
$a$ o $\ell$ no están acotados, y el costo de construirla, con solo el
factor $|Q|^{2\ell}$, lo es cuando $\ell$ no está acotado. El valor de
$\ell$ crece con la cantidad de variables de un mismo argumento del lado
izquierdo, pues por cada una queda libre el rango con que se la indexa.

En la modelación del capítulo~\ref{cap:modelacion-sat}, $G_{\mathrm{cons}}$
es la gramática de consistencia de una fórmula con $n$ ocurrencias de
variable booleana y $A$ su autómata booleano, con $|N| = O(n)$,
$|P| = O(n)$,
$|Q| = O(n^2)$ y $a = 1$. El número de predicados indexados es
polinomial,
\[
  |N'| = O(n\,(n^2)^2) = O(n^5).
\]
El número de cláusulas no lo es: la cláusula de $Eq_v$ para una variable
booleana con $r$ ocurrencias,
\[
  Eq_v(W_0\,b\,W_1\,b \cdots b\,W_r) \to L_{g_0}(W_0)\,L_{g_1}(W_1) \cdots L_{g_r}(W_r),
\]
tiene $r+1$ variables en su argumento del lado izquierdo, y $r$ puede ser
$O(n)$, de modo que $\ell = O(n)$. Entonces
\[
  |P'| = O\bigl(n\,(n^2)^{O(n)}\bigr) = 2^{O(n \log n)}.
\]
El tamaño de la gramática producto de $G_{\mathrm{cons}}$ y el autómata
booleano es
\[
  |G'| = |N'| + |T| + |P'| = O(n^5) + O(1) + 2^{O(n\log n)} = 2^{O(n\log n)}.
\]
Construirla cuesta lo mismo:
\[
  O(|P|\,\ell\,|Q|^{2\ell}) = O\bigl(n^2\,(n^2)^{O(n)}\bigr) = 2^{O(n\log n)}.
\]
Las cláusulas de $Eq_v$ son el término dominante.

En SATSRC el grado del polinomio con que se decide el vacío de la
intersección crece con la aridad de la sRCG, como se ve en la
subsección~\ref{subsec:consistencia-clase} del
capítulo~\ref{cap:rcg-sat}. La aridad $a = 1$ de $G_{\mathrm{cons}}$ se
debe a su no linealidad: un mismo argumento aparece en varios predicados
del lado derecho, y así la consistencia de cada variable booleana se
reconoce sin aumentar la aridad. A cambio se pierde la tratabilidad del
vacío de la intersección, polinomial con una sRCG y NP-completo con una
RCG no lineal.

Construir $G'$ sirve para decidir el vacío de la intersección, el
problema del capítulo~\ref{cap:vacio-interseccion-rcg-regular}. Como
$L(G') = L(G) \cap L(A)$, esa intersección es vacía si y solo si lo es
$L(G')$, esto es, si el predicado inicial $S'$ no es productivo. Para esa
decisión solo importan los predicados y cláusulas productivos, y muchos
de los indexados no lo son: no intervienen en el reconocimiento de
ninguna cadena de $L(G_{\mathrm{cons}}) \cap L(A)$, como se observó en la
sección~\ref{subsec:algoritmo}. En este trabajo no se da el costo de
decidir ese vacío, sino solo el objeto $G'$; limpiarlo de los predicados
y cláusulas no productivos y analizar ese costo quedan como problemas
abiertos.

En este capítulo se construyó la gramática producto $G'$ de una RCG $G$
y un autómata finito acíclico $A$, se demostró que
$L(G') = L(G) \cap L(A)$, se caracterizaron su tamaño y el costo de
construirla, exponenciales en el peor caso, y se mostró que muchos de
los predicados y cláusulas de $G'$ no son productivos. En el próximo
capítulo se presentan las conclusiones y recomendaciones del trabajo.

